Al giorno d'oggi lavorare con i dati provenienti da diverse sorgenti dati è diventata una necessità in qualsiasi contesto, ancora di più nell'ambito del Calcolo Scientifico. Per venire incontro a questo tipo di necessità Fly mette a disposizione degli sviluppatori delle funzionalità che permettono di accedere e utillizzare dati provenienti da sorgenti esterne. In questo momento è possibile in Fly manipolare dati presenti all'interno di file strutturati come CSV e fogli Excel grazie all'entità \verb|dataframe|. Essa permette di effettuare operazioni di lettura, scrittura e modifica su file CSV, fogli Excel e semplici txt. Inoltre, Fly supporta l'interazione con i database MySQL tramite le due entità \verb|sql| e \verb|query|: l'entità \verb|sql| permette la connessione e l'accesso a database MySQL sia locali che su Cloud, mentre l'entità \verb|query| permette di creare ed eseguire query per la manipolazione dei dati. \\
Negli ultimi anni alcuni ambiti hanno messo in secondo piano la grande affidabilità del modello relazionale. I database SQL di norma non sono molto scalabili e non puntano al massimo delle prestazioni. Questo risulta essere un grosso svantaggio per ambiti come il Calcolo Scientifico, dove \textbf{scalabilità} e \textbf{performance} diventano più importanti dell'affidabilità. La soluzione ricade nell'utilizzo dei database NoSQL. Grazie al dimensionamento orrizontalmente e ai diversi schemi flessibili utilizzabili per rappresentare i dati, i database NoSQL offrono una soluzione incentrata sul garantire la disponibilità dei dati. Questa soluzione risulta essere molto vantaggiosa per ambiti come il Calolo Scientifico dove vi è la necessità di lavorare con una mole di dati molto grande. È risultato quindi necessario implementare all’interno di Fly un supporto ai database NoSQL. Uno degli obiettivi di sviluppo principale è quello di rendere l'interazione con i database NoSQL all'interno di Fly omogenea. In questo modo lo sviluppatore potrà interagire con un cluster NoSQL distribuito su Cloud utilizzando la stessa sintassi necessaria per interagire con un cluster NoSQL distribuito in locale. Per poter realizzare questo obiettivo è risultato necessario implementare all'interno di Fly un supporto ai database NoSQL compatibili con MongoDB, in quanto al momento MongoDB è l'unico database NoSQL supportato dai diversi Cloud Provider. \\
In questo capitolo verrà descritta l'implementazione all'interno di Fly del supporto ai database NoSQL distribuiti sia in locale che su Cloud.

\section{Obiettivi}
Per permettere l'interazione con i database NoSQL all'interno di Fly sono stati posti alcuni obiettivi implementativi:
\begin{itemize}
    \item connettersi e accedere a cluster NoSQL sia locali che su Cloud;
    \item creare ed eseguire query per la manipolazione dei dati;
    \item fornire metodi per la gestione dei risultati;
    \item creare un meccanismo di interazione simile a quello già presente per i database MySQL.
\end{itemize}
Per poter interagire con i database MySQL Fly mette a disposizione dell'utente due entità:
\begin{itemize}
    \item \verb|sql| - entità che permette di stabilire una connessione ad un database SQL;
    \item \verb|query| - entità che permette di creare ed eseguire query sul database.
\end{itemize}
Grazie all'entità \verb|query| l'utente può effettuare diverse tipologie di query:
\begin{itemize}
    \item \verb|select|: query di selezione che restituisce come risultato una tabella;
    \item \verb|insert|: per aggiungere nuovi elementi all'interno di una tabella;
    \item \verb|update|: per modificare elementi presenti all'interno di una tabella;
    \item \verb|delete|: query di eliminazione che restituisce come risultato il numero di righe coinvolte.
\end{itemize}
Poiché uno degli obiettivi implementativi è proprio quello di creare un interazione con i database NoSQL il quanto più simile a quella già presenete per i database MySQL, l'utilizzo dell'entità \verb|query| verrà estesa ai database NoSQL. Inoltre, in modo analogo al funzionamento dell'entità \verb|sql|, verrà aggiunta una nuova entità \verb|nosql| per permettere all'utente di stabilire una connessione ad un database NoSQL.

\section{Implementazione}
\subsection{Connessione al database}
Come già introdotto nel capitolo precedente, la possibilità di stabilire una connessione ad un database NoSQL è stata resa possibile mediante l'implementazione dell'entità \verb|nosql|. L'obiettivo principale è quello di creare un meccanismo di connessione simile a quello già presente per i database SQL. La dichiarazione di una variabile di tipo \verb|sql| è composta dai seguenti parametri \cite{giuseppe}:
\begin{itemize}
    \item \textbf{db_name}: stringa rappresentante il nome del database;
    \item \textbf{resource_group} (solo in caso il database sia su Azure): stringa rappresentante il nome del ResourceGroup in cui è contenuta l’istanza del servizio DBaaS;
    \item \textbf{instance} (solo in caso il database sia su Cloud): stringa rappresentante il nome dell’istanza del servizio DBaaS su cui è in esecuzione il database;
    \item \textbf{user}: stringa rappresentante il nome utente per l’accesso al database;
    \item \textbf{password}: stringa rappresentante la password di accesso al database.
\end{itemize}
I database NoSQL, a differenza dei database SQL, partizionano il database in un insieme di collezioni. Per questo motivo l'entità \verb|nosql| deve permette di stabilire una connessione ad una precisa collezione del database. La dichiarazione di una variabile di tipo \verb|nosql| è composta quindi dai seguenti parametri:
\begin{itemize}
    \item \textbf{endpoint} (solo in caso non si usano i parametri \textbf{resource_group} e \textbf{instance}): stringa rappresentante l'endpoint del database;
    \item \textbf{resource_group} (solo in caso il database sia su Azure): stringa rappresentante il nome del ResourceGroup in cui è contenuta l’istanza del servizio DBaaS;
    \item \textbf{instance} (solo in caso il database sia su Azure): stringa rappresentante il nome dell’istanza del servizio DBaaS su cui è in esecuzione il database;
    \item \textbf{db_name}: stringa rappresentante il nome del database;
    \item \textbf{collection}: stringa rappresentante il nome della collezione del database;
    \item \textbf{properties} (opzionale): stringa rappresentante il path assoluto del file log4j.properties.
\end{itemize}
MongoDB utilizza \textbf{Log4J}, una libreria Java sviluppata dalla Apache Software Foundation che permette di mettere a punto un ottimo sistema di logging per tenere sotto controllo il comportamento di una applicazione, sia in fase di sviluppo che in fase di test e messa in opera del prodotto finale. Per poter quindi stabilire una connessione ad un database MongoDB è necessario specificare il path del file log4j.properties. Questo file contiene le proprietà che verranno utilizzate per il sistema di logging del cluster MongoDB. Questo parametro è però opzionale, infatti se l'utente decide di non assegnare nessun valore a questo parametro, Fly genera in automatico un file log4j.properties contenenti le seguenti proprietà:
\begin{lstlisting}[language=JAVA,caption={File log4j.proprietes generato in automatico da Fly.}, label={lst:log}]
log4j.appender.stdout.Target=System.out
log4j.appender.stdout.layout=org.apache.log4j.PatternLayout
log4j.appender.stdout=org.apache.log4j.ConsoleAppender
log4j.rootLogger=INFO, stdout
log4j.appender.stdout.layout.ConversionPattern=%d{yy/MM/dd HH\:mm\:ss} %p %c{2}\: %m%n
\end{lstlisting}  
Come possiamo notare l'oggetto \verb|nosql| segue una sintassi leggermente diversa dall'oggetto \verb|sql|. I database MongoDB utilizzano un URI per definire le connessioni tra applicazioni e istanze MongoDB. La stringa di connessione URI include i seguenti componenti:
\begin{itemize}
    \item \verb|mongodb://| - prefisso obliggatorio per identificare che si tratta di una stringa nel formato di connessione standard;
    \item \verb|username:password@| - se specificato, il client tenterà di autenticare l'utente nel \textbf{authSource} \cite{authSource};
    \item \verb|host[:port]| - l'host (e il numero di porta facoltativo) in cui è in esecuzione l'istanza MongoDB;
    \item \verb|?<options>| - opzionale. Una stringa di query che specifica le opzioni specifiche della connessione come \verb|<name>=<value>| coppie.
\end{itemize}
La stringa di connessione URI permette quindi di specificare le credenziali di accesso da utilizzare per stabilire una connessione ad un cluster MongoDB. Per questo motivo i parametri \verb|user| e \verb|password| non sono presenti all'interno dell'entità \verb|nosql|. Risulta invece necessario aggiungere il parametro \verb|collection| per poter specificare il nome della collezione a cui l'utente intende collegarsi. L'utilizzo della stringa di connessione URI permette all'utente di indicare l'endpoint del cluster MongoDB ovunque esso sia distribuito. Infatti, grazie a questa sintassi, l'utente può facilmente connettersi ad un database MongoDB distribuito in locale, su Azure oppure su AWS. Per quanto riguarda la connessione ad un database su Azure è possibile inoltre sostituire l'utilizzo del parametro \verb|client| con i parametri \verb|resourceGroup| e \verb|instance|. Questi due parametri consentono a Fly di ottenere in modo del tutto astratto all'utente l'endpoint del database NoSQL a cui si intende accedere. \\
Dopo la dichiarazione della varaibile, se il database è in esecuzione su Azure e si intende utilizzare i parametri \verb|resourceGroup| e \verb|instance|, occorre utilizzare la parola chiave \verb|on| seguita dalla variabile di \verb|type = "azure"|.
\begin{lstlisting}[language=FLY,caption={Dichiarazione di entità nosql per la connesione ad un database tramite l'utilizzo dell'endpoint.}, label={lst:nosql:1}]
var dbConnLocal = [type = "nosql", endpoint = "mongodb:127.0.0.1:27017", db_name = "nomeDatabase", collection = "nomeCollection"]
\end{lstlisting}  
\begin{lstlisting}[language=JAVA,caption={Dichiarazione di entità nosql per la connessione ad un database distribuito su Azure tramite l'utilizzo dei parametri resourceGroup e instance.}, label={lst:nosql:2}]
var dbConnLocal = [type = "nosql",  resource_group = "nomeResourceGroup ", instance = "nomeIstanza",  db_name = "nomeDatabase", collection = "nomeCollection"] on cloudAz
\end{lstlisting}  

\subsection{Creazione ed esecuzione di query}
Come già introdotto in precedenza, Fly mette a disposizione dell'utente l'entità \verb|query| per poter effettuare operazioni di inserimento, selezione, modifica e cancellazioni sui database MySQL. La dichiarazione di una variabile di tipo \verb|query| è composta, oltre dal parametro \verb|type="query"| che ne specifica appunto il tipo, dai seguenti parametri \cite{giuseppe}:
\begin{itemize}
    \item \textbf{query_type}: stringa rappresentante il tipo della query da eseguire (\verb|select|, \verb|insert|, \verb|update|, \verb|delete|);
    \item \textbf{connection}: nome della variabile di tipo 
